% Generated by scripts/assemble.py; corrections belong in source/corrections.json.
\examhead{31}{首都师范大学 2024 年数学分析试卷}{34}
% p035-b13; source PDF page 35
\par\Needspace{5\baselineskip}
\blockanchor{p035-b13}
\majorquestion{一}计算下列极限. 每题 8 分, 共 24 分.\par
% p035-b14; source PDF page 35
\blockanchor{p035-b14}
\question{1}$\lim_{x\to0^{+}}(\cos\sqrt{x})^{\frac{1}{x}}$\par
% p035-b15; source PDF page 35
\blockanchor{p035-b15}
\question{2}$\lim_{n\to \infty}a_n = a > 0,$ 求 $\lim_{n\to \infty}\sqrt[n]{a_n}$.\par
% p035-b16; source PDF page 35
\blockanchor{p035-b16}
\question{3}$\lim_{x\to 0}\frac{(1 + \tan x)^{\frac{1}{4}} + (1 - \sin x)^{\frac{1}{4}} - 2}{x^2}.$\par
% p035-b17; source PDF page 35
\par\Needspace{5\baselineskip}
\blockanchor{p035-b17}
\majorquestion{二}(10 分) 讨论 $f(x, y) = \left(1 + \frac{2}{x^2}\right)^{\frac{x^4}{x^2 + y^2}}$ 在 (0, 0) 处累次极限与重极限的存在性, 若存在, 求其值.\par
% p035-b18; source PDF page 35
\par\Needspace{5\baselineskip}
\blockanchor{p035-b18}
\majorquestion{三}(10 分) 证明: $f(x) = \ln x$ 在 $(1, +\infty)$ 上一致连续, 但在 $(0, 1)$ 上不一致连续.\par
% p035-b19; source PDF page 35
\par\Needspace{5\baselineskip}
\blockanchor{p035-b19}
\majorquestion{四}(10 分) 设 $f(x)$ 在 $[a, b]$ 上可导, 满足 $f(a) = f(b) = 0$, $f_{+}^{\prime}(a)f_{-}^{\prime}(b) > 0$, 证明: 至少存在不同的两点 $\xi, \eta \in (a, b)$, 使得 $f^{\prime}(\xi) = f^{\prime}(\eta) = 0$.\par
% p035-b20; source PDF page 35
\par\Needspace{5\baselineskip}
\blockanchor{p035-b20}
\majorquestion{五}(12 分) 证明: $f(x, y) = x^3 - 4x^2 + 2xy - y^2$ 在 $\mathbb{R}^2$ 上有唯一的极值点, 且该极值点是极大值点, 不是最大值点.\par
% p035-b21; source PDF page 35
\par\Needspace{5\baselineskip}
\par\Needspace{6\baselineskip}
\blockanchor{p035-b21}
\majorquestion{六}(14 分) 求三重积分\par
\nopagebreak[4]
% p035-b22; source PDF page 35
\blockanchor{p035-b22}
\displayformula{\iiint_ {\Omega} 2 \sqrt {x ^ {2} + y ^ {2}} d x d y d z.}
% p035-b23; source PDF page 35
\blockanchor{p035-b23}
其中 $\Omega$ 为柱面 $y = \sqrt{2x - x^2}$ 及平面 $z = 0, z = a (a > 0), y = 0$ 所围成的区域.\par
% p036-b00; source PDF page 36
\par\Needspace{5\baselineskip}
\blockanchor{p036-b00}
\majorquestion{七}(15 分) 定义 $f(x) = \begin{cases} \frac{1}{x} e^{-\frac{1}{x^2}}, & x \neq 0; \\ 0, & x = 0. \end{cases}$\par
% p036-b01; source PDF page 36
\blockanchor{p036-b01}
(1) (5 分) 证明: $f(x)$ 在点 $x = 0$ 处连续且可导.\par
% p036-b02; source PDF page 36
\blockanchor{p036-b02}
(2) (5 分) 证明: $f'(x)$ 在 $\mathbb{R}$ 上连续.\par
% p036-b03; source PDF page 36
\blockanchor{p036-b03}
(3) (5 分) 求 $f(x)$ 的单调区间, 最大值点, 最小值点.\par
% p036-b04; source PDF page 36
\par\Needspace{5\baselineskip}
\blockanchor{p036-b04}
\majorquestion{八}(15 分) 函数项级数 $\sum_{n=2}^{\infty} \frac{xe^{-nx}}{\ln n}$.\par
% p036-b05; source PDF page 36
\blockanchor{p036-b05}
(1) (7 分) 求此函数项级数的收敛区间.\par
% p036-b06; source PDF page 36
\blockanchor{p036-b06}
(2) (8 分) 设 $a > 0$, 证明此函数项级数在 $[a, +\infty)$ 上一致收敛.\par
% p036-b07; source PDF page 36
\par\Needspace{5\baselineskip}
\par\Needspace{6\baselineskip}
\blockanchor{p036-b07}
\majorquestion{九}（15分）设 $f_{1}(x)$ 在 $[a,b]$ 上可积， $A$ 为一给定的数，定义\par
\nopagebreak[4]
% p036-b08; source PDF page 36
\blockanchor{p036-b08}
\displayformula{f _ {n + 1} (x) = A + \int_ {a} ^ {x} f _ {n} (t) \mathrm{d} t.}
% p036-b09; source PDF page 36
\blockanchor{p036-b09}
其中  $x \in [a, b]$ ,  $n = 1, 2, \cdots$ .\par
% p036-b10; source PDF page 36
\blockanchor{p036-b10}
(1) (8 分) 证明: 函数列 $\{f_n(x)\}$ 在 $[a, b]$ 上一致收敛.\par
% p036-b11; source PDF page 36
\blockanchor{p036-b11}
(2) (7 分) 记 $\{f_n(x)\}$ 的极限函数为 $f(x)$, 证明: $f(x)$ 在 $[a, b]$ 上可微.\par
% p036-b12; source PDF page 36
\par\Needspace{5\baselineskip}
\par\Needspace{6\baselineskip}
\blockanchor{p036-b12}
\majorquestion{十}(15 分) 计算曲线积分\par
\nopagebreak[4]
% p036-b13; source PDF page 36
\blockanchor{p036-b13}
\displayformula{I = \int_ {L} x \ln (x ^ {2} + y ^ {2} - 1) \mathrm{d} x + y \ln (x ^ {2} + y ^ {2} - 1) \mathrm{d} y.}
% p036-b14; source PDF page 36
\blockanchor{p036-b14}
其中 L 是被积函数定义域内从点  $(2,0)$  到点  $(0,2)$  的逐段光滑曲线.\par
% p036-b15; source PDF page 36
\par\Needspace{5\baselineskip}
\blockanchor{p036-b15}
\majorquestion{十一}(10 分) 设 $f(x)$ 在 $[a, b]$ 上可积, 证明: $\lim_{\lambda \to +\infty} \int_{a}^{b} f(x) \sin \lambda x \, \mathrm{d}x = 0$.\par
